\subsubsection{Lax-Milgram 引理}\label{subsec:lax-milgram-lemma}

为建立椭圆方程的存在性理论而发展的第一个工具通常称为 Lax-Milgram 引理; 不过由于它的重要性, 将其列为定理. 这一结果只是将 Riesz 表示定理 (Theorem~\ref{thm:hilbert-riesz-representation}) 推广到不一定对称的双线性型.

\begin{Theorem}{Lax-Milgram}
\label{thm:lax-milgram}
设 $H$ 为 Hilbert 空间,
\begin{equation}\label{eq:lax-milgram-bilinear-mapping-domain}
B:H\times H\to\mathbb{R}
\end{equation}
为双线性映射. 假设存在正常数 $c_1$ 和 $c_2$, 使得
\begin{equation}\label{eq:lax-milgram-bilinear-boundedness}
|B[x,y]|\leq c_1\|x\|_H\|y\|_H\quad\text{for all }x\text{ and }y\text{ in }H
\end{equation}
且
\begin{equation}\label{eq:lax-milgram-bilinear-coercivity}
B[x,x]\geq c_2\|x\|_H^2\quad\text{for all }x\in H.
\end{equation}
那么对每个 $f\in H^*$, 都存在唯一的 $y\in H$, 使得
\begin{equation}\label{eq:lax-milgram-variational-solution}
B[x,y]=f(x)\quad\text{for all }x\in H.
\end{equation}
此外, 存在与 $f$ 无关的常数 $C$, 使得
\begin{equation}\label{eq:lax-milgram-solution-dual-norm-bound}
\|y\|_H\leq C\|f\|_{H^*}.
\end{equation}

\end{Theorem}

\begin{Remark}
\label{rem:coercive-form-energy-estimate}
若映射 $B$ 以某个 $c_2>0$ 满足 \eqref{eq:lax-milgram-bilinear-coercivity}, 则称其为\textbf{强制的}. 不等式 \eqref{eq:lax-milgram-bilinear-coercivity} 可看作能量估计. (该不等式说明, 能量 (范数的平方) 的增长速度最多与双线性型一样快.)
\end{Remark}

\begin{Remark}
\label{rem:lax-milgram-nonsymmetric-significance}
注意, 由 \eqref{eq:lax-milgram-bilinear-boundedness} 和 \eqref{eq:lax-milgram-bilinear-coercivity},
\begin{equation}\label{eq:coercive-bilinear-induced-equivalent-norm}
\|x\|_B:=\sqrt{B[x,x]}
\end{equation}
与 $H$ 上的原范数等价. 此外, 若 $B$ 对称, 即
\begin{equation}\label{eq:symmetric-bilinear-form-condition}
B[x,y]=B[y,x]\quad\text{for all }x\text{ and }y\text{ in }H,
\end{equation}
则 $B[x,y]$ 在 $H$ 上定义一个新的内积. 因而在这种情形下, Riesz 表示定理 (Theorem~\ref{thm:hilbert-riesz-representation}) 直接表明: 对每个 $f\in H^*$, 都存在唯一的 $y\in H$, 使得 \eqref{eq:lax-milgram-variational-solution} 成立. 因此, Lax-Milgram 引理的重要之处在于不要求 $B$ 对称.
\end{Remark}

现在证明 Lax-Milgram 引理.

\begin{Proof}
对每个固定的 $y\in H$, 映射
\begin{equation}\label{eq:lax-milgram-fixed-second-argument-functional}
H\ni x\mapsto B[x,y]\in\mathbb{R}
\end{equation}
有界且线性, 即为 $H^*$ 中的元素. 因而由 Riesz 表示定理 (Theorem~\ref{thm:hilbert-riesz-representation}), 存在唯一的 $z\in H$, 使得
\begin{equation}\label{eq:lax-milgram-riesz-representing-vector}
B[x,y]=(x,z)\quad\text{for all }x\in H.
\end{equation}
由于对每个固定的 $y\in H$ 都能得到唯一的 $z\in H$, 可以用下式定义映射 $A:H\to H$:
\begin{equation}\label{eq:lax-milgram-associated-operator-definition}
z=:A(y).
\end{equation}
现在, \eqref{eq:lax-milgram-variational-solution} 的解的存在性问题转化为 $A$ 的可逆性问题. 即对任意 $f\in H^*$, 设 $z\in H$ 为满足下式的唯一元素:
\begin{equation}\label{eq:lax-milgram-data-riesz-representative}
(x,z)=f(x)\quad\text{for all }x\in H.
\end{equation}
那么, 若对每个 $z\in H$ 均能找到 $A(y)=z$ 的唯一解 $y\in H$, 则 $y$ 为下式的唯一解:
\begin{equation}\label{eq:lax-milgram-operator-equation-equivalence}
B[x,y]=(x,A(y))=(x,z)=f(x)\quad\text{for all }x\in H.
\end{equation}
现在注意 $A$ 的一些基本性质.\begin{enumerate}\item $A$ 线性. 为此注意\[\begin{aligned}
(x,A(\alpha_1y_1+\alpha_2y_2))&:=B[x,\alpha_1y_1+\alpha_2y_2]\\
&=\alpha_1B[x,y_1]+\alpha_2B[x,y_2]\\
&=\alpha_1(x,A(y_1))+\alpha_2(x,A(y_2))\\
&=(x,\alpha_1A(y_1)+\alpha_2A(y_2)).
\end{aligned}\]由于这对任意 $x$, $\alpha_i$ 和 $y_i$ 成立, 线性得证.\item $A$ 有界. 利用 \eqref{eq:lax-milgram-bilinear-boundedness}, 得到
\begin{equation}\label{eq:lax-milgram-associated-operator-norm-calculation}
\|A(y)\|^2=(A(y),A(y))=B[A(y),y]\leq c_1\|A(y)\|\|y\|.
\end{equation}
约去公因子, 得到
\begin{equation}\label{eq:lax-milgram-associated-operator-upper-bound}
\|A(y)\|\leq c_1\|y\|.
\end{equation}
\item $A$ 的值域在 $H$ 中稠密. 为此, 由 \eqref{eq:lax-milgram-bilinear-coercivity} 注意到, 若 $y\in\mathcal{R}(A)^\perp$, 则
\begin{equation}\label{eq:lax-milgram-associated-operator-range-dense}
c_2\|y\|^2\leq B[y,y]=(y,A(y))=0.
\end{equation}
\item $A$ 有下界. 再次利用 \eqref{eq:lax-milgram-bilinear-coercivity} (这次对任意 $y\in H$), 得到
\begin{equation}\label{eq:lax-milgram-associated-operator-lower-calculation}
c_2\|y\|^2\leq B[y,y]=(y,A(y))\leq\|y\|\|A(y)\|
\end{equation}
或
\begin{equation}\label{eq:lax-milgram-associated-operator-bounded-below}
\|A(y)\|\geq c_2\|y\|.
\end{equation}
\item 将此与 Problem~\ref{ex:closed-bounded-below-operator-closed-range} 结合, 说明 $A$ 的值域 $\mathcal{R}(A)$ 是闭的. 由于 $\mathcal{R}(A)$ 稠密, $A$ 为满射.\end{enumerate}于是 $A$ 可逆, 从而得到唯一解 $y$ 的存在性. 最后, 估计 \eqref{eq:lax-milgram-solution-dual-norm-bound} 由 Riesz 表示定理 (Theorem~\ref{thm:hilbert-riesz-representation}) 和 $A$ 有下界这一事实得到.
\end{Proof}
